% ============================================================================
% Fichier  : partie4/P04_C20_audit_smsi.tex
% Titre    : Auditer un Système de Management de la Sécurité de l'Information
% Auteur   : MINKA MI NGUIDJOÏ Thierry Emmanuel
% Version  : 1.0 -- novembre 2026
% ============================================================================
% Base : ACSI C16 (trois audits de sécurité ; chaîne du SMSI ; « une mesure
% déclarée n'est pas une mesure en place »).
% Sources : ISO/IEC 27001:2022 (clauses 4 à 10, annexe A, déclaration
% d'applicabilité) ; ISO/IEC 27002:2022 (93 mesures en 4 thèmes) ;
% ISO/IEC 27007:2020 (audit d'un SMSI) ; ISO 19011:2018 ; ISO/IEC 27006
% (organismes de certification, édition à vérifier) ; loi n° 2010/012,
% art. 13 (audit de sécurité obligatoire), texte officiel lu le 06/11/2026.
% Faits nouveaux (SANAGA) : périmètre « centre de données de Douala » ;
% déclaration d'applicabilité 93 mesures, 81 retenues, 12 exclues dont 3
% non justifiées ; lecture de la chaîne du SMSI 7/2/3.
% Corrigés : corriges/Corr_C20.tex.
% ============================================================================

\chapter{Auditer un Système de Management de la Sécurité}
\label{chap:audit-smsi}
\chaptermark{Auditer un SMSI}

\epigraph{%
    « Un SMSI ne promet pas qu'il n'y aura pas d'incident.\\
    Il promet qu'on saura pourquoi, et qu'on corrigera. »%
}{Principe du management de la sécurité}

\niveauChapitre{Intermédiaire}{%
    Maîtrise des Chapitres~\ref{chap:gestion-risques} (appréciation des
    risques), \ref{chap:normes} (ISO/IEC~27007) et
    \ref{chap:referentiels} (ISO/IEC~27002)%
}

\begin{boxobjectifs}
\begin{itemize}
    \item Distinguer trois audits de sécurité~: l'audit légal, l'audit d'un
          SMSI et l'audit technique.
    \item Lire l'ISO/IEC~27001 clause par clause et savoir quelle preuve
          demander pour chacune.
    \item Auditer une déclaration d'applicabilité et la justification des
          exclusions.
    \item Conduire un audit interne de SMSI selon l'ISO~19011 et
          l'ISO/IEC~27007, et distinguer les étapes d'une certification.
    \item Coter les maillons de la chaîne du SMSI et lire le résultat.
    \item Rédiger un écart d'audit de SMSI (non-conformité majeure ou
          mineure) et le relier à l'observation en sept éléments.
\end{itemize}
\end{boxobjectifs}

\bigskip

% ============================================================================
\section*{Question directrice : la banque gère-t-elle sa sécurité, ou la subit-elle~?}
% ============================================================================

À Douala, la BANQUE SANAGA prépare à la fois l'audit de sécurité annuel de
l'ANTIC et la certification ISO/IEC~27001 de son centre de données. Sa
directrice générale a reçu du RSSI un classeur de 300~pages~: politique,
analyse de risques, déclaration d'applicabilité, procédures. Elle demande à
Mme~EKAMBI, directrice de l'audit interne~: «~Sommes-nous prêts~?~»

Mme~EKAMBI sait qu'un classeur ne prouve rien~: elle doit vérifier que
chaque document est appliqué, et que la banque sait corriger ce qui ne
l'est pas.

\medskip
\textbf{Plan du chapitre.}
\begin{enumerate}
    \item Trois audits de sécurité (section~\ref{sec:c20-trois}).
    \item L'ISO/IEC~27001, clause par clause (section~\ref{sec:c20-clauses}).
    \item La déclaration d'applicabilité (section~\ref{sec:c20-dda}).
    \item Conduire l'audit d'un SMSI (section~\ref{sec:c20-conduire}).
    \item Coter la chaîne du SMSI de la BANQUE SANAGA
          (section~\ref{sec:c20-sanaga}).
    \item Rédiger un écart (section~\ref{sec:c20-ecart}).
\end{enumerate}

\separateur

% ============================================================================
\section{Trois Audits de Sécurité}
\label{sec:c20-trois}
% ============================================================================

\begin{tableaucours}{L{2.4cm}L{3.2cm}L{3.4cm}L{3.2cm}}{%
    \tableauHeader{Audit} &
    \tableauHeader{Fondement} &
    \tableauHeader{Objet} &
    \tableauHeader{Destinataire}
}
Audit de sécurité légal & Loi n°~2010/012, art.~13 et~32 & Politique,
mesures, procédures et moyens, au regard des risques & Autorité
compétente~; rapport confidentiel \\
Audit d'un SMSI & ISO/IEC~27001 (§~9.2) et 27007~; certification par un
organisme accrédité & Le système de management~: risques, mesures,
amélioration & Direction~; organisme de certification \\
Audit technique & Lettre de mission et note d'autorisation & Configurations,
vulnérabilités, flux, comptes & RSSI, DSI \\
\end{tableaucours}
\vspace{-0.8em}
\begin{center}{\small\itshape Tableau~20.1. Trois audits de sécurité}\end{center}

Les trois se complètent. L'audit d'un SMSI vérifie que l'organisation
\emph{gère} sa sécurité~; l'audit technique (Chapitres~\ref{chap:audit-technique}
et~\ref{chap:intrusion}) vérifie que les mesures sont \emph{effectivement
en place}~; l'audit légal couvre les deux.

\begin{boxjuris}
\textbf{\loicam{n°~2010/012}, article~13, alinéa~1} (texte officiel)~:
\textit{« Sont soumis à un audit de sécurité obligatoire, les réseaux de
communications électroniques et les systèmes d'information, des
opérateurs, des autorités de certification et des fournisseurs de
services de communications électroniques. »}

\medskip
\textbf{Article~32, alinéas~1 et~2} (texte officiel)~: \textit{« Les réseaux
de communications électroniques et les systèmes d'information sont soumis
à un régime d'audit de sécurité obligatoire et périodique de leurs
systèmes de sécurité par l'Agence. L'audit de sécurité et les mesures
d'impact de gravité sont effectués chaque année ou lorsque les
circonstances l'exigent. »}

\medskip
\textbf{Analyse.} L'article~13 désigne des catégories d'acteurs~;
l'article~32 vise plus largement les systèmes d'information, avec un
audit annuel conduit par l'Agence, dont les modalités sont fixées par voie
réglementaire (texte d'application à vérifier). Pour une banque s'ajoute
la réglementation COBAC (Chapitre~\ref{chap:droit}). L'auditeur vérifie le
fondement exact avant de qualifier un audit d'« obligatoire ».
\end{boxjuris}

% ============================================================================
\section{L'ISO/IEC 27001, Clause par Clause}
\label{sec:c20-clauses}
% ============================================================================

\subsection{La structure de la norme}

L'ISO/IEC~27001 suit la structure commune des normes de systèmes de
management~: sept clauses d'exigences (4 à 10), et une annexe~A qui
énumère les 93~mesures de sécurité de l'ISO/IEC~27002:2022~\cite{iso27001_2022,iso27002_2022}.

\begin{tableaucours}{cL{3.2cm}L{4.4cm}L{4.0cm}}{%
    \tableauHeader{Clause} &
    \tableauHeader{Thème} &
    \tableauHeader{Exigence principale (reformulée)} &
    \tableauHeader{Preuve à demander}
}
4 & Contexte & Comprendre l'organisation, les parties intéressées~;
définir le périmètre & Déclaration de périmètre, analyse du contexte \\
5 & Leadership & Engagement de la direction~; politique~; rôles &
Politique signée, nominations \\
6 & Planification & Apprécier et traiter les risques~; objectifs de
sécurité & Méthode et registre des risques, plan de traitement,
déclaration d'applicabilité \\
7 & Support & Ressources, compétences, sensibilisation, communication,
documentation & Plan de formation, maîtrise des documents \\
8 & Fonctionnement & Mettre en œuvre le traitement, réapprécier les risques
& Preuves d'exécution des mesures, mises à jour du registre \\
9 & Évaluation & Surveiller, mesurer, auditer (§~9.2), revue de direction
(§~9.3) & Indicateurs, programme et rapports d'audit interne, compte
rendu de revue \\
10 & Amélioration & Traiter les non-conformités, améliorer en continu &
Registre des non-conformités et actions correctives \\
\end{tableaucours}
\vspace{-0.8em}
\begin{center}{\small\itshape Tableau~20.2. Les clauses 4 à 10 de
l'ISO/IEC~27001:2022 et leurs preuves}\end{center}

\subsection{La chaîne du SMSI}

Pour l'auditeur, le SMSI est une chaîne, et chaque maillon laisse une
pièce. Un maillon manquant rompt la chaîne~: des mesures sans appréciation
des risques sont arbitraires~; une appréciation sans mesures est un
exercice de papier~; des mesures sans évaluation ne s'améliorent pas.

\begin{tableaucours}{L{3.6cm}L{4.4cm}L{4.2cm}}{%
    \tableauHeader{Maillon} &
    \tableauHeader{Pièce} &
    \tableauHeader{Question d'audit}
}
Périmètre & Déclaration de périmètre & Les systèmes critiques y
figurent-ils~? \\
Appréciation des risques & Registre des risques de sécurité & Est-elle
datée, revue, cohérente avec les actifs~? \\
Traitement & Plan de traitement, déclaration d'applicabilité & Chaque
exclusion de mesure est-elle justifiée~? \\
Mesures & Procédures, configurations, journaux & Chaque mesure retenue
fonctionne-t-elle~? \\
Évaluation & Rapports d'audit interne, revue de direction & Les écarts
sont-ils suivis jusqu'à leur clôture~? \\
\end{tableaucours}
\vspace{-0.8em}
\begin{center}{\small\itshape Tableau~20.3. La chaîne du SMSI et ses
preuves}\end{center}

\begin{boxinfo}
\textbf{Une mesure déclarée n'est pas une mesure en place.} La déclaration
d'applicabilité dit qu'une mesure est retenue~; elle ne dit pas qu'elle
fonctionne. Sans preuve de fonctionnement, l'auditeur cote \NEK{}, et le
taux d'incertitude du SMSI mesure la distance entre le papier et la
réalité.
\end{boxinfo}

% ============================================================================
\section{La Déclaration d'Applicabilité}
\label{sec:c20-dda}
% ============================================================================

\subsection{Ce qu'elle contient}

\begin{boxdef}
\termecle{Déclaration d'applicabilité}\index{déclaration
d'applicabilité}~: document exigé par l'ISO/IEC~27001 qui liste les mesures
de l'annexe~A, indique pour chacune si elle est retenue ou exclue, justifie
ce choix, et précise si la mesure est mise en œuvre.
\end{boxdef}

L'auditeur en fait trois lectures~: \emph{complétude} (les 93~mesures
sont-elles toutes traitées~?), \emph{justification} (chaque exclusion
s'appuie-t-elle sur l'appréciation des risques ou sur l'absence de
l'activité concernée~?) et \emph{réalité} (chaque mesure dite « mise en
œuvre » l'est-elle~?).

\subsection{La déclaration de la BANQUE SANAGA}

\begin{boxscenario}{Opération SANAGA. Douze exclusions}
\textbf{Faits.} La déclaration d'applicabilité du centre de données de
Douala traite les 93~mesures~: 81 sont retenues, 12 exclues. Mme~EKAMBI
examine les douze justifications.

\medskip
\begin{tableaucours}{L{5.2cm}L{4.2cm}c}{%
    \tableauHeader{Exclusion (exemples)} &
    \tableauHeader{Justification donnée} &
    \tableauHeader{Justifiée~?}
}
Travail à distance & Aucun télétravail autorisé & \OUI{} \\
Sécurité du développement & « La banque ne développe pas » & \NON{} \\
Test de sécurité dans le développement & Idem & \NON{} \\
Sécurité des services en nuage & « Aucun service en nuage » & \NON{} \\
Huit autres exclusions & Activité absente, vérifiée & \OUI{} \\
\end{tableaucours}

\medskip
\textbf{Constatation.} Trois exclusions sur douze ne sont pas justifiées~:
DIGICORE AFRIQUE paramètre et modifie le système central (activité de
développement), et la messagerie de la banque est un service en nuage
souscrit en 2025.

\medskip
\textbf{Analyse} (confiance élevée). Les trois exclusions portent
exactement sur les zones où le risque s'est réalisé (changements du
système central) ou peut se réaliser (données dans le nuage). Une
exclusion injustifiée est un trou dans la chaîne, au maillon
« traitement ».
\end{boxscenario}

\begin{boxalert}
\textbf{Violation typique.} Exclure une mesure parce qu'elle est difficile à
mettre en œuvre.

\smallskip
\textbf{Règle~:} une mesure ne s'exclut que si le risque qu'elle traite
n'existe pas dans le périmètre, ou s'il est accepté par une décision
écrite de la direction, fondée sur l'appréciation des risques. La
difficulté n'est pas une justification.
\end{boxalert}

% ============================================================================
\section{Conduire l'Audit d'un SMSI}
\label{sec:c20-conduire}
% ============================================================================

\subsection{Audit interne, audit de certification}

\begin{tableaucours}{L{3.0cm}L{4.4cm}L{4.8cm}}{%
    \tableauHeader{} &
    \tableauHeader{Audit interne (§~9.2)} &
    \tableauHeader{Audit de certification}
}
Par qui & L'organisation, ou un prestataire pour son compte & Un organisme
de certification accrédité (ISO/IEC~27006) \\
Pour quoi & Vérifier la conformité et l'efficacité, alimenter la revue de
direction & Délivrer, maintenir ou renouveler un certificat \\
Rythme & Programme planifié, au moins sur le cycle du certificat & Étape~1
(revue documentaire), étape~2 (évaluation sur site), puis audits de
surveillance et de renouvellement \\
Résultat & Constats et actions correctives & Non-conformités majeures ou
mineures, décision de certification \\
\end{tableaucours}
\vspace{-0.8em}
\begin{center}{\small\itshape Tableau~20.4. Audit interne et audit de
certification d'un SMSI}\end{center}

L'ISO/IEC~27007 complète l'ISO~19011 pour les SMSI~: elle précise ce qu'il
faut vérifier pour chaque clause et comment échantillonner les mesures de
l'annexe~A~\cite{iso27007_2020,iso19011_2018}.

\subsection{Échantillonner les mesures}

Tester les 81~mesures retenues dans une seule mission est rarement
possible. L'auditeur les échantillonne selon le risque, en s'assurant que
le programme couvre toutes les mesures sur le cycle de trois ans du
certificat.

\begin{boxprocedure}
\textbf{Programme d'audit interne d'un SMSI sur trois ans}
\begin{enumerate}
    \item Chaque année~: clauses 4 à 10 (le système), mesures liées aux
          risques élevés du registre, mesures qui ont connu un incident.
    \item Sur trois ans~: toutes les mesures retenues de la déclaration
          d'applicabilité, réparties par thème.
    \item À chaque changement majeur (nouveau système, nouveau prestataire)~:
          audit ciblé des mesures concernées.
    \item Pour chaque mesure testée~: une fiche d'assertion, une preuve de
          fonctionnement sur la période, une cote.
\end{enumerate}
\end{boxprocedure}

\subsection{Interroger la direction}

La clause~5 (leadership) et la clause~9.3 (revue de direction) se
vérifient auprès de la direction, pas du RSSI. L'auditeur demande le
compte rendu de la dernière revue de direction et y cherche les éléments
d'entrée exigés (résultats d'audit, incidents, objectifs, risques) et les
décisions prises.

\begin{boxmetier}{RSSI}
\textbf{Préparer l'audit d'un SMSI}
\begin{enumerate}
    \item Tenir la déclaration d'applicabilité à jour à chaque changement
          de périmètre (nouveau service en nuage, nouveau prestataire).
    \item Conserver, pour chaque mesure retenue, une preuve de
          fonctionnement datée.
    \item Présenter à la revue de direction les résultats des audits et
          l'état des actions correctives.
    \item Traiter les écarts de l'audit interne avant l'audit de
          certification, et pouvoir le montrer.
\end{enumerate}
\end{boxmetier}

% ============================================================================
\section{La Chaîne du SMSI de la BANQUE SANAGA}
\label{sec:c20-sanaga}
% ============================================================================

\begin{boxlecture}{Le SMSI du centre de données}
Mme~EKAMBI cote douze contrôles répartis sur les cinq maillons de la
chaîne.

\medskip
\begin{tableaucours}{L{3.2cm}L{6.2cm}c}{%
    \tableauHeader{Maillon} &
    \tableauHeader{Contrôle} &
    \tableauHeader{Cote}
}
Périmètre & Le périmètre est écrit et approuvé & \OUI{} \\
 & Les systèmes critiques y figurent (système central) & \OUI{} \\
Risques & La méthode d'appréciation est écrite & \OUI{} \\
 & Le registre est revu depuis la migration de 2026 & \NON{} \\
Traitement & La déclaration d'applicabilité traite les 93~mesures &
\OUI{} \\
 & Les exclusions sont justifiées & \NON{} \\
Mesures & Les mesures d'accès fonctionnent (échantillon) & \OUI{} \\
 & Les mesures de journalisation fonctionnent & \NEK{} \\
 & Les mesures de gestion des changements fonctionnent & \NEK{} \\
Évaluation & Un audit interne du SMSI a eu lieu dans l'année & \OUI{} \\
 & Ses écarts sont suivis jusqu'à clôture & \NEK{} \\
 & Une revue de direction a eu lieu dans l'année & \OUI{} \\
\end{tableaucours}

\medskip
\textbf{Constatation.} 7~\OUI{}, 2~\NON{}, 3~\NEK{}~: $\mu \approx 0{,}58$~;
$U = 0{,}25$~; $H \approx 1{,}38$~bit.

\medskip
\textbf{Analyse} (confiance moyenne). Conformité opaque~: le système est
bien construit sur le papier, mais un quart des contrôles, tous dans les
maillons « mesures » et « évaluation », ne peut pas être établi. Les deux
\NON{} (registre non revu, exclusions injustifiées) sont liés~: la
migration n'a pas été répercutée dans le SMSI. La réponse à la directrice
générale est~: \emph{pas encore prêts}, et l'on sait où agir.

\medskip
\hfill\lecture{0{,}58}{0{,}25}{1{,}38}
\end{boxlecture}

% ============================================================================
\section{Rédiger un Écart}
\label{sec:c20-ecart}
% ============================================================================

\subsection{Majeure ou mineure}

Dans un audit de certification, les écarts se classent en deux niveaux.

\begin{boxdef}
\termecle{Non-conformité majeure}\index{non-conformité majeure}~: absence
ou défaillance d'une exigence qui remet en cause la capacité du système à
atteindre ses résultats (par exemple, aucune appréciation des risques, ou
aucun audit interne). Elle empêche la certification tant qu'elle n'est pas
corrigée.

\medskip
\termecle{Non-conformité mineure}~: défaillance ponctuelle qui ne remet pas
en cause le système (une procédure non mise à jour, un compte rendu non
signé).
\end{boxdef}

La gradation ressemble à celle des observations du cours
(Chapitre~\ref{chap:conclure}), mais sa règle est fixée par l'organisme de
certification. L'auditeur interne qui prépare une certification adopte la
même grille, pour que la direction sache ce qui bloquera.

\subsection{Un écart de la BANQUE SANAGA}

\begin{boxobjet}{Exclusions injustifiées}
\textbf{Norme~:} ISO/IEC~27001:2022, clause~6.1.3 (déclaration
d'applicabilité, justification des exclusions).

\textbf{Constatation~:} trois des douze exclusions ne sont pas justifiées~:
sécurité du développement, test de sécurité dans le développement, sécurité
des services en nuage~; or DIGICORE AFRIQUE paramètre le système central et
la messagerie est hébergée en nuage depuis 2025 (contrat DIGICORE~;
contrat de messagerie).

\textbf{Risques~:} des changements du système central sans exigences de
sécurité~; des données de messagerie hors du périmètre de maîtrise.

\textbf{Classement~:} non-conformité majeure si l'organisme considère que
le traitement des risques est défaillant sur des activités critiques~;
mineure sinon. L'auditrice recommande de la traiter comme majeure.

\textbf{Recommandations~:} réviser la déclaration d'applicabilité~;
réapprécier les risques après la migration~; intégrer DIGICORE et la
messagerie au périmètre des mesures.
\end{boxobjet}

\begin{boxmetier}{Auditeur de sécurité}
\textbf{Relier l'audit du SMSI et l'audit technique.} L'auditeur de
sécurité qui teste les configurations (Chapitre~\ref{chap:audit-technique})
transmet ses constats au maillon « mesures » du SMSI~: une mesure
déclarée en place et trouvée défaillante en test est un \NON{} dans les
deux audits. Les deux audits partagent les preuves, et se renforcent.
\end{boxmetier}

\begin{boxmetier}{Responsable de la conformité}
\textbf{Ne pas confondre certificat et conformité.} Un certificat
ISO/IEC~27001 atteste que le SMSI du périmètre déclaré respecte la norme à
la date de l'audit. Il ne dit pas que la banque respecte le règlement
COBAC, ni la loi sur les données personnelles. Le responsable de la
conformité tient son propre registre des obligations
(Chapitre~\ref{chap:audit-conformite}).
\end{boxmetier}

\begin{boxinfo}
\textbf{Ce que dit la recherche.} Le cours lit la chaîne du SMSI avec
$\mu$, $U$ et $H$~\cite{minka2026nek}. L'intérêt est de distinguer un SMSI
dont les mesures échouent (\NON{} nombreux) d'un SMSI dont on ne sait pas si
elles fonctionnent (\NEK{} nombreux)~: le second est plus fréquent dans les
organisations qui préparent une certification, et c'est lui qui expose aux
non-conformités majeures.
\end{boxinfo}

\separateur

\begin{boxresume}
\textbf{Ce qu'il faut retenir de ce chapitre~:}
\begin{itemize}
    \item \textbf{Trois audits de sécurité}~: légal, d'un SMSI, technique~;
          ils se complètent et partagent leurs preuves.
    \item \textbf{L'ISO/IEC~27001} a sept clauses d'exigences (4 à 10) et une
          annexe~A de 93~mesures~; chaque clause a sa preuve.
    \item \textbf{Le SMSI est une chaîne}~: périmètre, risques, traitement,
          mesures, évaluation~; une mesure déclarée sans preuve de
          fonctionnement est \NEK{}.
    \item \textbf{La déclaration d'applicabilité} s'audite en complétude,
          justification et réalité~; la difficulté n'est pas une
          justification.
    \item \textbf{L'audit interne} échantillonne les mesures selon le risque,
          et les couvre toutes sur trois ans~; la certification procède en
          deux étapes, puis par surveillance.
    \item \textbf{Un écart} est majeur s'il remet en cause le système,
          mineur sinon.
\end{itemize}
\end{boxresume}

\begin{boxinfo}
\textbf{Points de contrôle de l'élève.} Je sais~: distinguer les trois
audits de sécurité~; citer les clauses 4 à 10 et une preuve pour
chacune~; auditer une exclusion~; construire un programme d'audit sur
trois ans~; coter la chaîne du SMSI~; classer un écart.
\end{boxinfo}

\bigskip

\begin{boxexo}[title={\ding{45}\quad Exercice~20.1. Quel audit~? \niveaubase}]
Pour chaque demande, indiquez l'audit de sécurité correspondant.
\begin{enumerate}[(a)]
    \item L'ANTIC annonce son passage annuel chez un opérateur de
          téléphonie.
    \item Le RSSI veut savoir si les règles des pare-feu de Garoua
          correspondent à la matrice de flux.
    \item La BANQUE SANAGA veut un certificat pour son centre de données.
    \item Le comité d'audit demande si les écarts du dernier audit interne
          du SMSI ont été clos.
\end{enumerate}
\end{boxexo}

\begin{boxexo}[title={\ding{45}\quad Exercice~20.2. Quelle clause~? \niveaubase}]
Rattachez à une clause de l'ISO/IEC~27001~: (a) la politique de sécurité
n'est pas signée~; (b) aucune revue de direction depuis deux ans~; (c) le
registre des risques n'a pas été mis à jour après la migration~; (d) les
agents n'ont reçu aucune sensibilisation~; (e) une non-conformité de
l'audit précédent n'a pas d'action corrective.
\end{boxexo}

\begin{boxexo}[title={\ding{45}\quad Exercice~20.3. Auditer une exclusion \niveaumoyen}]
LiZbiZ prépare son propre SMSI et exclut la mesure sur la sécurité des
fournisseurs, « car nos fournisseurs ne touchent pas au SI ». Or un
prestataire installe et exploite les centres de Garoua et de Bafoussam.
Analysez, cotez et rédigez la recommandation.
\end{boxexo}

\begin{boxexo}[title={\ding{45}\quad Exercice~20.4. Le programme sur trois ans \niveaumoyen}]
Proposez le programme d'audit interne du SMSI de la BANQUE SANAGA pour
2026, 2027 et 2028~: ce qui est audité chaque année, ce qui est réparti, ce
qui est déclenché par un événement.
\end{boxexo}

\begin{boxexo}[title={\ding{45}\quad Exercice~20.5. Recalculer la lecture \niveaumoyen}]
Recalculez $\mu$, $U$ et $H$ de la chaîne du SMSI de la BANQUE SANAGA (posez
le calcul). Après un mois de travail, les trois \NEK{} deviennent deux
\OUI{} et un \NON{}. Recalculez, situez dans le quadrant et commentez.
\end{boxexo}

\begin{boxexo}[title={\ding{45}\quad Exercice~20.6. Majeure ou mineure~? \niveaumoyen}]
Classez et justifiez~: (a) aucun audit interne depuis la certification
initiale~; (b) une procédure de sauvegarde datée de 2021~; (c) la
déclaration d'applicabilité ne traite que 80 des 93~mesures~; (d) le
compte rendu d'une revue de direction n'est pas signé.
\end{boxexo}

\begin{boxexo}[title={\ding{45}\quad Exercice~20.7. Le certificat d'un prestataire \niveaumoyen}]
DIGICORE AFRIQUE présente un certificat ISO/IEC~27001 dont le périmètre
est « hébergement ». Que vérifiez-vous sur le certificat (organisme,
accréditation, dates, périmètre, déclaration d'applicabilité)~? Que vaut-il
pour l'exploitation du système central~?
\end{boxexo}

\begin{boxexo}[title={\ding{45}\quad Exercice~20.8. Interroger la direction \niveaumoyen}]
Rédigez cinq questions pour l'entretien de Mme~EKAMBI avec la directrice
générale sur la clause~5 et la revue de direction, et la pièce attendue
pour chaque réponse.
\end{boxexo}

\begin{boxexo}[title={\ding{45}\quad Exercice~20.9. Le SMSI après la migration \niveauexpert}]
Rédigez l'observation en sept éléments relative au registre des risques
non revu depuis la migration, et expliquez comment elle s'articule avec
l'écart sur les exclusions.
\end{boxexo}

\begin{boxexo}[title={\ding{45}\quad Exercice~20.10. Prêts pour la certification~? \niveauexpert}]
Rédigez la note de Mme~EKAMBI à la directrice générale (une page)~: état
de la chaîne, écarts qui bloqueraient une certification, plan d'action et
calendrier réaliste avant l'étape~1.
\end{boxexo}

\bigskip

\begin{boxinfo}
\textbf{Pour aller plus loin.}
\begin{itemize}
    \item ISO/IEC~27001:2022~\cite{iso27001_2022} et
          ISO/IEC~27002:2022~\cite{iso27002_2022}.
    \item ISO/IEC~27007:2020~\cite{iso27007_2020}~: l'audit d'un SMSI.
    \item ISO~19011:2018~\cite{iso19011_2018}~: le programme d'audit.
\end{itemize}
\end{boxinfo}

\begin{boxinfo}
\textbf{Transition.} L'audit du SMSI dit si la banque gère sa sécurité~;
il reste à vérifier que les mesures techniques sont réellement en place.
Le Chapitre~\ref{chap:audit-technique} audite l'infrastructure~:
segmentation, durcissement, protocoles, vulnérabilités et journalisation.
\end{boxinfo}
